% 为加快编译速度，撰写报告时插入的所有图片都以草稿形式显示，报告写完后把第10行 "\PassOptionsToPackage{draft}{graphicx}" 注释掉，即可恢复图片显示。
% 实验报告封面通过插入.pdf文件的方式引入，封面的.pdf文件可用Word模板生成。
% 如在正文中\indent 缩进无效，请尝试写两次\indent，或使用\qquad 缩进。
% 在tcolorbox环境中浮动体机制会失效，在正文中插入浮动体时必须加入[H]参数，强制浮动体在代码位置插入。如插入图片浮动体应使用\begin{figure}[H]，否则会报错。这相当于放弃了浮动功% % 能，此时的浮动体实质上可以用\begin{minipage}{} 替代，也可以直接用本模板中提供的透明盒子环境替代它。
% 本模板会在实验报告的首页底部、尾页顶部、中间页顶部和底部自动补画框线，有时会出现章节分隔线和页面框线重复的问题。此时请删除对应章节末尾画出章节分隔线的指令，并人为在章节末尾插入一% 个透明盒子环境，环境中写一个小垂直空格，如\vspace{1pt}，将下一章节的标题挤到下一页去，即可手动解决双框线问题。

% =========================== 导言区开始 =================================
\documentclass[a4paper, zihao=5, UTF8, AutoFakeBold]{ctexart}
\PassOptionsToPackage{draft}{graphicx}

%% ===== 参考文献 =====
\usepackage[backend=biber, style=gb7714-2015]{biblatex}
\addbibresource{references.bib}

%% ===== 排版宏包 =====
\usepackage{pdfpages}
\usepackage{fancyhdr}
\usepackage{zhnumber}
\usepackage{fontspec}
\usepackage{setspace}
\usepackage{titlesec}
\usepackage[most]{tcolorbox}
\usepackage{tabularray}
\usepackage{booktabs}
\usepackage{geometry}
\usepackage{caption}
\usepackage{float}
\usepackage{dblfloatfix}
\usepackage{siunitx}
\usepackage{mathtools}
\usepackage{unicode-math}
\usepackage{derivative}
\usepackage{xfrac}
\usepackage{physics2}
\usephysicsmodule{ab,braket,diagmat,doubleprod}
\usepackage{romannum}
\usepackage{graphicx}
\usepackage{xcolor}
\usepackage{diagbox}
\usepackage{adjustbox}
\usepackage[nameinlink]{cleveref}
\usepackage{etoolbox}
\usepackage{tikz}

%% ===== 大框样式 =====
\tcbset
{
    reportbox/.style=
    {
        breakable, enhanced,
        colback  = white,
        colframe = black,
        boxrule  = 0.6pt,
        arc      = 0pt,
        left = 2pt, right = 3pt,
        top  = 2pt, bottom = 2pt,
        boxsep   = 0pt,
        overlay first={
            \draw[black, line width=0.6pt]
                (frame.south west) -- (frame.south east);
        },
        overlay middle={
            \draw[black, line width=0.6pt]
                (frame.north west) -- (frame.north east);
            \draw[black, line width=0.6pt]
                (frame.south west) -- (frame.south east);
        },
        overlay last={
            \draw[black, line width=0.6pt]
                (frame.north west) -- (frame.north east);
        },
    },
    segmentation style=
    {
        solid,
        draw=black,
        line width=0.6pt
    }
}

%% ===== 正文内容区环境 =====
\newenvironment{contentarea}
{
    \par\vspace{4pt}
    \leftskip=\cpad\rightskip=\cpad
    \setlength{\parindent}{2em}
    \par\vspace{4pt}
}

%% ===== 内容区左右内边距 =====
\newlength{\cpad}
\setlength{\cpad}{5pt}

%% ===== 透明盒子 =====
\newtcolorbox{transparentbox}
{
    enhanced,
    frame hidden, interior hidden,
    boxrule=0pt,
    boxsep=0pt,
    left=0pt, right=0pt,
    top=0pt, bottom=0pt,
    center,
}

%% ===== 页面设置 =====
\geometry{a4paper,top=2.5cm, bottom=1.5cm, left=2cm, right=2cm}

\pagestyle{fancy}
\fancyhf{}                       
\fancyfoot[r]{\thepage}          
\renewcommand{\headrulewidth}{0pt}
\renewcommand{\footrulewidth}{0pt}

%% ===== 行距 段距 栏距 =====
\setstretch{1.25}
\setlength{\parskip}{0pt}
\setlength{\parindent}{2em}
\setlength{\columnsep}{2.5em}

%% ===== 字体设置 =====
\setmainfont{Times New Roman}
\setCJKmainfont{SimSun}
\setmathfont{STIX Two Math}

%% ===== 图头、表头格式 =====
\captionsetup{font=small,format=hang,justification=centering}

%% ===== 节标题格式 =====
\titleformat{\section}[block]{\hspace{5pt}\bfseries\LARGE}{\thesection .\hspace{-4pt}}{0.5em}{}
\titleformat{\subsection}[block]{\hspace{5pt}\bfseries\Large}{\thesubsection .\hspace{-4pt}}{0.5em}{}
\titleformat{\subsubsection}[block]{\hspace{5pt}\bfseries\large}{\thesubsubsection .\hspace{-4pt}}{0.5em}{}

\renewcommand{\thesection}{\arabic{section}}
\renewcommand{\thesubsection}{\arabic{section}.\arabic{subsection}}
\renewcommand{\thesubsubsection}{\arabic{section}.\arabic{subsection}.\arabic{subsubsection}}

\titlespacing*{\section}{0pt}{0pt}{0pt}
\titlespacing*{\subsection}{0pt}{0pt}{0pt}
\titlespacing*{\subsubsection}{0pt}{0pt}{0pt}

%% ===== 自定义节标题 =====
\newcommand{\secheader}[1]{\noindent\hspace{\cpad}{\bfseries #1}}

%% ===== 计数器 =====
\newcounter{thepart}

%% ===== 引用格式 =====
% 公式编号
\numberwithin{equation}{thepart}
\renewcommand{\theequation}{\arabic{thepart}.\arabic{equation}}
\crefformat{equation}{#2式~(#1)#3}
\Crefformat{equation}{#2式~(#1)#3}

% 图编号
\numberwithin{figure}{thepart}
\renewcommand{\thefigure}{\arabic{thepart}.\arabic{figure}}
\crefformat{figure}{#2图~#1#3}
\Crefformat{figure}{#2图~#1#3}

% 表编号
\numberwithin{table}{thepart}
\renewcommand{\thetable}{\arabic{thepart}.\arabic{table}}
\crefformat{table}{#2表~#1#3}
\Crefformat{table}{#2表~#1#3}

% 章节编号
\crefformat{section}{#2第#1节#3}
\Crefformat{section}{#2第#1节#3}

% ===== 记号 =====
% 矢量标记
\newcommand{\vb}[1]{\boldsymbol{#1}}
\newcommand{\va}[1]{\overrightarrow{#1}}

% 微分
\newcommand{\ff}{\delta}
\newcommand{\dd}{\mathrm{d}}

% 算符
\newcommand{\grad}{\nabla}
\newcommand{\diver}{\nabla\cdot}
\newcommand{\curl}{\nabla\times}
\newcommand{\lapla}{\nabla^2}
\newcommand{\dAlem}{\Box}

% 数值与单位
\sisetup{range-phrase = {\text{～}}}
\DeclareSIUnit{\angstrom}{\text{\AA}}

% =========================== 内容区开始 =================================
\begin{document}

%% =====  封面 =====
%\includepdf{}

\begin{tcolorbox}[reportbox]

%% ===== 一、实验目的 =====
\setcounter{thepart}{1}
\vspace{1pt}
\secheader{一、实验目的}
\vspace{-8pt}
\begin{contentarea}

\noindent 示例内容\\示例内容\\示例内容\\示例内容\\示例内容\\示例内容\\示例内容\\示例内容\\示例内容\\示例内容\\
示例内容\\示例内容\\示例内容\\示例内容\\示例内容\\示例内容\\示例内容\\示例内容\\示例内容\\示例内容\\

\end{contentarea}
\tcbline

%% ===== 二、实验原理 =====
\setcounter{thepart}{2}
\setcounter{section}{0}
\setcounter{figure}{0}
\setcounter{table}{0}
\vspace{-2pt}
\secheader{二、实验原理}
\vspace{-8pt}
\begin{contentarea}

\noindent 示例内容\\示例内容\\示例内容\\示例内容\\示例内容\\示例内容\\示例内容\\示例内容\\示例内容\\示例内容\\
示例内容\\示例内容\\示例内容\\示例内容\\示例内容\\示例内容\\示例内容\\示例内容\\示例内容\\示例内容\\

\end{contentarea}
\tcbline
 
%% ===== 三、实验仪器 =====
\setcounter{thepart}{3}
\setcounter{section}{0}
\setcounter{figure}{0}
\setcounter{table}{0}
\vspace{-2pt}
\secheader{三、实验仪器}
\vspace{-8pt}
\begin{contentarea}

\noindent 示例内容\\示例内容\\示例内容\\示例内容\\示例内容\\示例内容\\示例内容\\示例内容\\示例内容\\示例内容\\
示例内容\\示例内容\\示例内容\\示例内容\\示例内容\\示例内容\\示例内容\\示例内容\\示例内容\\示例内容\\

\end{contentarea}
\tcbline

%% ===== 四、实验内容与步骤 =====
\setcounter{thepart}{4}
\setcounter{section}{0}
\setcounter{figure}{0}
\setcounter{table}{0}
\vspace{-2pt}
\secheader{四、实验内容与步骤}
\vspace{-8pt}
\begin{contentarea}

\noindent 示例内容\\示例内容\\示例内容\\示例内容\\示例内容\\示例内容\\示例内容\\示例内容\\示例内容\\示例内容\\
示例内容\\示例内容\\示例内容\\示例内容\\示例内容\\示例内容\\示例内容\\示例内容\\示例内容\\示例内容\\

\end{contentarea}
\tcbline

%% ===== 五、数据处理与分析 =====
\setcounter{thepart}{5}
\setcounter{section}{0}
\setcounter{figure}{0}
\setcounter{table}{0}
\vspace{-2pt}
\secheader{五、数据处理与分析}
\vspace{-8pt}
\begin{contentarea}

\noindent 示例内容\\示例内容\\示例内容\\示例内容\\示例内容\\示例内容\\示例内容\\示例内容\\示例内容\\示例内容\\
示例内容\\示例内容\\示例内容\\示例内容\\示例内容\\示例内容\\示例内容\\示例内容\\示例内容\\示例内容\\

\end{contentarea}
\tcbline

%% ===== 六、结果陈述 =====
\setcounter{thepart}{6}
\setcounter{section}{0}
\setcounter{figure}{0}
\setcounter{table}{0}
\vspace{-2pt}
\secheader{六、结果陈述}
\vspace{-8pt}
\begin{contentarea}

\noindent 示例内容\\示例内容\\示例内容\\示例内容\\示例内容\\示例内容\\示例内容\\示例内容\\示例内容\\示例内容\\
示例内容\\示例内容\\示例内容\\示例内容\\示例内容\\示例内容\\示例内容\\示例内容\\示例内容\\示例内容\\

\end{contentarea}
\tcbline

%% ===== 七、思考题 =====
\setcounter{thepart}{7}
\setcounter{section}{0}
\setcounter{figure}{0}
\setcounter{table}{0}
\vspace{-2pt}
\secheader{七、思考题}
\vspace{-8pt}
\begin{contentarea}

\noindent 示例内容\\示例内容\\示例内容\\示例内容\\示例内容\\示例内容\\示例内容\\示例内容\\示例内容\\示例内容\\
示例内容\\示例内容\\示例内容\\示例内容\\示例内容\\示例内容\\示例内容\\示例内容\\示例内容\\示例内容\\

\end{contentarea}
\tcbline

% ===== 指导教师批阅意见 =====
\vspace{-2pt}
\secheader{指导教师批阅意见}
\vspace{3cm}
\tcbline

% ===== 成绩评定 =====
\vspace{-2pt}
\secheader{成绩评定}
\vspace{-8pt}
\begin{table}[H]
\centering
\begin{tblr}{|Q[c,m]|Q[c,m]|Q[c,m]|Q[c,m]|Q[c,m]|Q[c,m]|}
    \hline
    {预习\\（20分）} & {操作及记录\\（40分）} & {数据处理与分析\\（20分）} & {结果陈述与讨论\\（10分）} & {思考题\\（10分）} & 总分 \\
    \hline
    \\
    \\
    \hline
\end{tblr}
\end{table}

\vspace{1pt}
\end{tcolorbox}

\end{document}